% Generated draft: requires mathematical and source review.
\documentclass[11pt]{article}
\usepackage[margin=25mm]{geometry}
\usepackage{fontspec}
\setmainfont{Latin Modern Roman}
\usepackage{amsmath,amssymb,mathtools,unicode-math}
\setmathfont{Latin Modern Math}
\usepackage{enumitem,needspace,xurl,hyperref}
\setlength{\parindent}{0pt}
\setlength{\parskip}{5pt}
\setlength{\emergencystretch}{3em}
\newcommand{\N}{\mathbb{N}}
\newcommand{\Z}{\mathbb{Z}}
\newcommand{\Q}{\mathbb{Q}}
\newcommand{\R}{\mathbb{R}}
\newcommand{\C}{\mathbb{C}}
\newcommand{\F}{\mathbb{F}}
\newcommand{\A}{\mathbb{A}}
\newcommand{\PP}{\mathbb{P}}
\newcommand{\E}{\mathbb{E}}
\newcommand{\eps}{\varepsilon}
\newcommand{\dd}{\,\mathrm{d}}
\begin{document}
\section*{UnsolvedMath ID: 30000890}
\textbf{Source problem code:} catalogue entry
\subsection{Definitions and mathematical statement}
The source poses an SYZ mirror-symmetry question for a manifold $X$ with a singular special Lagrangian torus fibration $f:X\to B$. To explain the usual terms, in a complex $d$-dimensional K\"ahler setting with symplectic form $\omega$ and holomorphic volume form $\Omega$, a real $d$-dimensional submanifold $L$ is special Lagrangian of phase $\theta$ if $\omega|_L=0$ and $\operatorname{Im}(e^{-i\theta}\Omega)|_L=0$, with the compatible calibrated orientation. The extract does not specify the full ambient Calabi--Yau or relative/logarithmic data. Let $B^0\subseteq B$ be the regular-value locus and write $L_b=f^{-1}(b)\simeq(S^1)^d$ there. The usual dual torus at $b$ is \[\widehat L_b=\operatorname{Hom}(H_1(L_b;\Z),U(1)),\] the space of rank-one unitary local systems on $L_b$. The question asks whether the mirror is the dual torus fibration over the same base $B$, equipped with a holomorphic map $W:\check X\to\C$ whose superpotential appropriately counts holomorphic discs in $X$ with boundary on the original fibers $L_b$.\par This is the source's proposed picture, not an assertion that the uncorrected dual bundle over $B^0$ already extends over all singular fibers or automatically has the required complex structure. The extract leaves unspecified the singular compactification, disc-count definition and weights, correction/gluing procedure, and criterion for being a mirror. It also does not distinguish a compact Calabi--Yau mirror with possibly zero potential from a setting requiring a Landau--Ginzburg potential; no nonzero $W$ or particular Maslov-index restriction is added. The 2026-08-21 triage labels the general question open, reports rigorous semi-flat, toric, and other special realizations, and says existence/control of singular special-Lagrangian torus fibrations and the full construction remain open in general. The source question begins with a given fibration; that separate existence issue is retained as a literature qualification, not substituted for its hypothesis.
\subsection{Short English statement}
Given a singular special Lagrangian torus fibration, ask whether its mirror comes from dualizing the fibers over the same base and adding a holomorphic disc-counting superpotential. The singular extension, disc counts, corrections, and precise meaning of mirror are not specified.
\subsection{Sources}
\begin{itemize}\item Local UnsolvedMath registry, permanent ID 30000890, source problem catalogue entry, \emph{Strominger--Yau--Zaslow Mirror-Symmetry Conjecture}, proposed year 2007. The original labels the conjecture SYZ; no individual proposer metadata is supplied.\item Supplied citation: \emph{Tropical Geometry} (2008), Oberwolfach Reports, DOI 10.4171/owr/2007/57, pp.\ 3303--3368: \url{https://doi.org/10.4171/owr/2007/57}. Supplied literature label: original Oberwolfach report.\item Supplied literature label: special-Lagrangian fibrations and SYZ, \url{https://arxiv.org/abs/math/0011179}.\item Supplied literature label: SYZ mirror symmetry for toric Calabi--Yau manifolds, \url{https://arxiv.org/abs/1006.3830}.\end{itemize}
\end{document}
